% ============================================================
% 全国大学生数学建模竞赛（国赛）中文论文模板 · XeLaTeX
% ------------------------------------------------------------
% 这份模板只给你「结构 + 空位」，不给你任何可粘贴的正文章句。
% 每一节的标题旁边注明了它对应左边的哪个阶段：把任务卡里填过的
% 内容搬进来，才是这篇论文你自己的东西。
%
% 编译方式：xelatex 两遍（本应用点「编译」即自动完成）。
% 图形直接用沙箱产物文件名，例如 \includegraphics{curve.png}。
% 承诺页与编号页按赛区要求单独装订，不进这个文件。
% ============================================================
\documentclass[12pt,a4paper,UTF8]{ctexart}

% 国赛格式规范要求数字与西文用 Times New Roman。字体探测不到就退回默认，
% 别让一台没装该字体的机器直接编译失败。
\IfFontExistsTF{Times New Roman}{\setmainfont{Times New Roman}}{}
\usepackage{amsmath,amssymb}
\usepackage{geometry}
\geometry{a4paper,top=2.5cm,bottom=2.5cm,left=2.8cm,right=2.8cm,headheight=14pt}
\usepackage{graphicx}
\usepackage{booktabs}
\usepackage{tabularx}
\usepackage{longtable}
\usepackage{multirow}
\usepackage{array}
\usepackage{float}
\usepackage{caption}
\usepackage{enumitem}
\usepackage{setspace}
\usepackage{fancyhdr}
\usepackage{xcolor}
\usepackage[colorlinks=true,linkcolor=black,citecolor=black,urlcolor=black]{hyperref}

\graphicspath{{./}}

\captionsetup{font=small,labelsep=quad}
\captionsetup[table]{skip=4pt}
\renewcommand{\arraystretch}{1.25}

\pagestyle{fancy}
\fancyhf{}
\fancyfoot[C]{\zihao{5}\thepage}
\renewcommand{\headrulewidth}{0pt}

% 正文行距与段落间距：评阅人翻得快，太挤直接影响印象分
\linespread{1.35}
\setlength{\parskip}{0.4ex}

% ============================================================
% 题目与摘要页
% 题目要落到方法上，别只写"……的建模"。摘要页是评阅人唯一保证会
% 逐字看的部分，对应阶段 9。
% ============================================================
\begin{document}

\begin{center}
  {\zihao{-3}\heiti 这里填论文题目}\\[6pt]
\end{center}

\begin{center}
  {\zihao{4}\heiti 摘\quad 要}
\end{center}

% 摘要三段式：问题 → 方法（具体到模型名与求解器）→ 结论（带数值）。
% 逐问都要覆盖，缺哪一问就是丢哪一问的分。
这里写摘要第一段：赛题要解决的核心问题是什么，你的总体思路是什么。

这里写摘要第二段：针对问题一/二/三分别用了什么方法，模型的关键假设与求解算法叫什么，用什么工具求解、规模多大。

这里写摘要第三段：给出主要结果的数值（购电量、费用、误差、灵敏度幅度等），并说明结论在什么条件下成立。

\begin{flushleft}
  \zihao{-4}\textbf{关键词：} 关键词一；关键词二；关键词三；关键词四；关键词五
\end{flushleft}

\newpage

% ============================================================
\section{问题重述}
% 对应阶段 1「读题与问题拆解」。把任务卡「逐问拆解」「附件与数据」
% 两栏的内容改写成完整句子放进来，不要复制赛题原文整段。
% ============================================================

\subsection{问题背景}
问题背景用自己的话复述：研究对象是什么、要调控/优化/评价的对象是什么、为什么难。

\subsection{需要解决的问题}
问题一：……（已知条件、要求的结果、结果要交到哪里）

问题二：……

\subsection{附件与数据}
\begin{table}[H]
  \centering
  \caption{附件数据说明}
  \begin{tabularx}{\textwidth}{lXl}
    \toprule
    附件 & 内容与字段 & 规模 \\
    \midrule
    附件1 & 待填 & 待填 \\
    附件2 & 待填 & 待填 \\
    \bottomrule
  \end{tabularx}
\end{table}

% ============================================================
\section{问题分析}
% 对应阶段 4「模型选型」。这一节要写出「比较」而不是「宣布」：
% 每个小问有哪几条可选路线、各自的适用前提与代价、为什么最终选它。
% ============================================================

各小问的求解对象、数据形态与目标函数形式分别是什么，据此说明为什么问题一适合先做确定性问题、问题三为什么必须引入随机性或滚动调整。

% ============================================================
\section{模型假设}
% 对应阶段 3。每条假设一行，并在括号里注明它在后面哪一节被用到。
% 「假设数据真实可靠」这种空话不要写。
% ============================================================
\begin{enumerate}[leftmargin=2em,itemsep=2pt]
  \item 假设……（用于第 X 节的目标函数构造）
  \item 假设……（若不为真，影响是……）
\end{enumerate}

% ============================================================
\section{符号说明}
% 对应阶段 3「符号表」。符号全文唯一、必须带单位。
% ============================================================
\begin{longtable}{clc}
  \caption{主要符号及含义} \\
  \toprule
  符号 & 含义 & 单位 \\
  \midrule
  \endfirsthead
  \toprule
  符号 & 含义 & 单位 \\
  \midrule
  \endhead
  $x$ & 待填 & 待填 \\
  $C$ & 待填 & 元 \\
  \bottomrule
\end{longtable}

% ============================================================
\section{数据的预处理与探索分析}
% 对应阶段 2。写清缺失/异常怎么判定、怎么处理、为什么，并放至少
% 一张能支撑后续建模选择的分布图。图无坐标轴名称与单位＝白画。
% ============================================================

\begin{figure}[H]
  \centering
  % 把沙箱跑出来的图直接按文件名引用
  % \includegraphics[width=0.72\textwidth]{explore.png}
  \caption{这里写图题：它回答哪个问题}
  \label{fig:explore}
\end{figure}

\subsection{数据清洗规则}
清洗规则：缺失如何判定、如何处理；异常如何判定；量纲与时间粒度如何统一。

\subsection{探索性结论}
从数据里读出的、真正影响建模决定的两三条事实。

% ============================================================
\section{模型的建立与求解}
% 对应阶段 5 与 6。每个小问一节：先给决策变量与目标函数，再给约束
% （逐条标注来源：题目给定还是假设引入），最后写求解算法与结果。
% 公式要有编号并在正文里被引用，结果要有量纲解释。
% ============================================================

\subsection{问题一的模型}
决策变量、目标函数与约束。

\begin{equation}
  \min \; Z = \sum_{t=1}^{T} p_t \, x_t
  \label{eq:objective}
\end{equation}
其中 $p_t$ 为时段 $t$ 的电价，$x_t$ 为待填。式~\eqref{eq:objective} 的约束为：
\begin{equation}
  \begin{cases}
    \text{约束一} \\
    \text{约束二}
  \end{cases}
\end{equation}

\subsection{问题一的求解与结果}
求解思路、算法、规模与耗时；结果按赛题要求的表格格式给出。

\begin{table}[H]
  \centering
  \caption{结果表（按赛题指定格式填写）}
  \begin{tabular}{cc}
    \toprule
    量 & 数值 \\
    \midrule
    待填 & \\
    \bottomrule
  \end{tabular}
\end{table}

\subsection{问题二的模型}
与问题一的差异在哪里，为什么原模型不能直接用。

\subsection{问题二的求解与结果}
结果与解释，异常值怎么处理。

% ============================================================
\section{灵敏度分析}
% 对应阶段 7，这是评分表里的稳健性项，也是最常被跳过的一节。
% 扰动幅度要说明依据，结论要给出「在多大范围内不变、在哪里翻转」。
% ============================================================

\begin{table}[H]
  \centering
  \caption{参数扰动对目标值的影响}
  \begin{tabular}{ccc}
    \toprule
    参数 & 扰动幅度 & 目标值变化率 \\
    \midrule
    待填 & $\pm 10\%$ & 待填 \\
    \bottomrule
  \end{tabular}
\end{table}

% ============================================================
\section{模型的评价与改进}
% 优点、缺点都要具体到可核验的点上；改进方向要能接着做，不写
% "可以进一步推广"这类话。
% ============================================================

\subsection{模型优点}

\subsection{模型缺点与改进方向}

% ============================================================
\begin{spacing}{1.0}
\begin{thebibliography}{99}
% 参考文献只列你真正引用过的；正文里必须有 \cite 与之对应。
% 数量与格式不一致是提交前自查清单里的常见失分项。
  \bibitem{ref1} 作者. 书名[M]. 出版地: 出版社, 年.
  \bibitem{ref2} 作者. 题名[J]. 期刊名, 年, 卷(期): 起止页码.
\end{thebibliography}
\end{spacing}

% ============================================================
\section*{附录}
\addcontentsline{toc}{section}{附录}
% 附录：支撑代码（含可复现的运行说明）、补充图表、超出页数的推导。
% 源程序要能跑出正文里的结果，否则不要放。
% ============================================================
源代码清单、补充数据表、以及正文放不下但影响复现的材料。

\end{document}
